\section{New-Build Gentrification in Central Shanghai: Demographic Changes and Socioeconomic Implications}

\begin{description}
  \item[Authors:] He, S. (Shenjing He)
  \item[Year:] 2010 (published online 2 July 2009)
  \item[Journal:] Population, Space and Place, 16(5), 345--361
  \item[DOI:] 10.1002/psp.548
\end{description}

\subsection*{Domain Classification}

\begin{itemize}
  \item \textbf{Primary domain:} Geography (critical geography of gentrification and displacement)
  \item \textbf{Secondary domain(s):} Regional Science (multi-scale spatial analysis via Location Quotient indices); Urban Economics (land-lease reform and re-emergence of a land price gradient, treated only descriptively)
  \item \textbf{Keywords:} new-build gentrification; state-led/state-sponsored gentrification; exclusionary and direct displacement; Location Quotient (LQ) index; occupational class remaking; land-lease reform; socio-spatial polarization; neoliberal urbanism; sub-district (jiedao) spatial unit; comparative case-study design
\end{itemize}

\subsection*{Research Question}

The paper asks how new-build gentrification has reshaped the socioeconomic and occupational composition of central Shanghai's population, and what socioeconomic consequences this process has for displaced low-income households. Specifically, He examines (1) how the spatial distribution of ``potential gentrifiers'' and ``potential displacees'' changed in central Shanghai between 1990 and 2000, and (2) how the socioeconomic profiles of residents in redeveloped (gentrified) neighborhoods compare with those in adjacent, not-yet-redeveloped (pre-gentrification) neighborhoods, including the anticipated impacts of displacement on displacees' quality of life and socioeconomic prospects.

\subsection*{Theoretical Framework}

The paper is situated in the critical, political-economy tradition of gentrification research rather than in a formal economic equilibrium framework. It draws explicitly on the concept of \emph{new-build gentrification} (Davidson and Lees, 2005) --- gentrification manifested through demolition--rebuild redevelopment and large-scale luxury apartment construction rather than incremental rehabilitation of existing housing stock --- and on Smith's (1996) framing of gentrification as the ``class remake'' of urban landscapes. He situates Shanghai's experience within the debate on gentrification as a manifestation of globalization and neoliberal urbanism (Smith, 2002; Atkinson and Bridge, 2005; Lees, Slater and Wyly, 2008), while engaging Wu's (2008) counter-argument that China's neoliberalization is not simply an import of globalization but a ``great transformation'' resolving the crisis of capital accumulation under state socialism. Building on her own prior work on \emph{state-sponsored gentrification} under market transition (He, 2007), the paper's central theoretical contribution is to insist that Chinese new-build gentrification is driven by strong, direct state intervention --- in land supply, infrastructure investment, and resettlement housing provision --- rather than by market forces or global capital flows alone. Methodologically, the occupational classification of ``gentrifier proxy'' (Class 1) versus ``displacee proxy'' (Class 2) groups follows Atkinson (2000). Notably, the paper does not invoke or operationalize Neil Smith's (1979) rent gap model explicitly, despite the relevance of land-value dynamics to its empirical narrative; the theoretical engagement remains at the level of critical urban political economy (the ``eviction of critical perspectives'' debate: Slater, 2006, 2008; Wacquant, 2008; Watt, 2008; Smith, 2008) rather than formal land-market modeling.

\subsection*{Data}

\begin{itemize}
  \item \textbf{Source(s):} Chinese population census data (1953, 1964, 1982, 1990, and 2000 rounds); 2008 Shanghai Statistics Yearbook; a structured household survey conducted by the author in 2004 (500 questionnaires); 16 semi-structured in-depth interviews conducted during 2004 fieldwork; municipal administrative records on residential relocation and housing demolition (1995--2007).
  \item \textbf{Geographic scope:} Central Shanghai municipality, People's Republic of China --- specifically the nine central urban districts (Huangpu, Luwan, Xuhui, Changning, Jing'an, Hongkou, Putuo, Zhabei, Yangpu), with comparisons to Pudong New District and outskirt districts (e.g., Minhang, Baoshan, Jiading).
  \item \textbf{Spatial unit:} District level (1990 occupational data) and sub-district (\emph{jiedao}) level (2000 occupational data) for the Location Quotient analysis; individual housing estate/neighborhood level for the four case-study sites (Liangwancheng, Shenjiazhai, Xintiandi, Jing'anli).
  \item \textbf{Time period:} Population redistribution analyzed 1953--2007; detailed occupational-group analysis 1990--2000; industrial-sector employment analyzed 1990, 2000, and 2007; household survey and interviews conducted in 2004.
  \item \textbf{Sample size:} Metropolitan Shanghai population of 15.57 million (2000 census); household survey N = 500 (200 in Liangwancheng, 100 in Shenjiazhai, 100 in Xintiandi, 100 in Jing'anli); 16 in-depth interviews with residents, relocated households, and a residents'-committee officer.
  \item \textbf{Key variables:} Occupational classification collapsed into Class 1 (``gentrifier proxy'': managers and senior officers, professionals, associate professional and administrative staff) and Class 2 (``displacee proxy'': secretaries/sales/customer service, processing/plant/machine operatives, agricultural workers, others/unemployed); the Location Quotient (LQ) index of occupational concentration; industrial-sector employment counts; household-level socioeconomic indicators (age, educational attainment, occupation, working sector, annual family income, housing tenure category, second-home ownership); residents' self-predicted displacement impacts (access to public facilities, unemployment risk, social-network fragmentation, distance from city center).
  \item \textbf{Data limitations:} The 2000 census occupational data were released only as a 10\% sample and rescaled by a factor of 10 to approximate comparability with the 1990 full count, introducing potential small-area sampling noise. Sub-district-level occupational data are unavailable for 1990, forcing a coarser district-level comparison for that year. The four survey sites were selected purposively (matched on pre-existing built environment and development history, not randomly), and sample fractions of total neighborhood households vary widely (roughly 3\% in Liangwancheng, 7\%--8\% in the other three sites). Most importantly, the Xintiandi sample consists of consumers and visitors rather than former or current residents, because the area was fully converted to commercial/recreational use; the author acknowledges this comparability problem explicitly but retains the comparison on the grounds that ``new-build gentrification'' is conceptualized as class remaking of urban space rather than narrowly as residential upgrading.
\end{itemize}

\subsection*{Methodology}

\begin{itemize}
  \item \textbf{Empirical strategy:} A descriptive, comparative mixed-methods design. Quantitatively, He computes Location Quotient (LQ) indices from census occupational and industrial-sector data to track the spatial concentration of ``potential gentrifier'' and ``potential displacee'' populations across the metropolitan area, central area, individual districts, and sub-districts (1990 vs.\ 2000), supplemented by administrative statistics on residential relocation and housing demolition (1995--2007) and by industrial-sector employment change (1990, 2000, 2007). This is combined with a comparative case-study design contrasting two matched neighborhood pairs --- a pre-gentrification site and an adjacent post-gentrification (redeveloped) site --- via a structured household survey (Table 7) and 16 semi-structured in-depth interviews with displaced and relocated residents.
  \item \textbf{Identification:} The paper does not employ a causal identification strategy (no instrumental variable, regression discontinuity, difference-in-differences, or matching estimator). Comparisons are purely descriptive contrasts between purposively matched neighborhood pairs and between census years; there is no counterfactual estimation of what would have happened to residents or neighborhoods absent redevelopment.
  \item \textbf{Spatial treatment:} Spatial concentration is measured only through Location Quotient indices and choropleth mapping (Figure 2) at multiple nested scales (metropolitan, central area, district, sub-district, and concentric distance rings around People's Square at 2.5, 5, 7.5, and 10\,km). No formal spatial econometric model (e.g., spatial lag, spatial error, or Durbin specification), spatial weights matrix, or spatial autocorrelation statistic (e.g., Moran's I) is estimated; spatial dependence is visualized descriptively rather than modeled statistically.
  \item \textbf{Robustness checks:} The author triangulates across multiple data vintages (five census rounds spanning 1953--2000, plus 2008 Yearbook data) and across data types (household survey, in-depth interviews, and multiple informant categories --- inner-city residents, gentrified-area residents, and relocated suburban residents, including a residents'-committee officer). The comparability limitation of the Xintiandi sample is explicitly flagged rather than concealed.
  \item \textbf{Methodological innovation:} The paper adapts Atkinson's (2000) occupational gentrifier/displacee proxy and the LQ concentration index --- techniques developed in Western, market-economy contexts --- to a state-socialist transitional economy that lacks official statistics tracking gentrification or displacement directly, and where the very term ``gentrification'' is largely absent from official and even academic Chinese-language discourse.
\end{itemize}

\subsection*{Key Findings}

\begin{enumerate}
  \item \textbf{Central-city population decline}: Between 1990 and 2000, population in three central urban districts (Huangpu, Jing'an, Luwan) fell by 691,700, a decline of 36.4\%, while population in six other districts (including Pudong, Minhang, Baoshan, Jiading) rose by 2,998,400, an increase of 63.5\%. Within a 2.5\,km and 5\,km radius of People's Square, population fell by 20\% and 12\% respectively, while it rose by 8\% and 50\% within 7.5\,km and 10\,km radii --- a pattern the author argues reflects state-sponsored class displacement rather than conventional automobile-driven suburbanization.
  \item \textbf{Scale of relocation and demolition}: Between 1995 and 2007, 946,424 households were relocated and 51.052 million\,m\textsuperscript{2} of housing was demolished in Shanghai, peaking at 98,714 relocated households in 2002 (Table 2).
  \item \textbf{Occupational class remaking (LQ evidence)}: Central Shanghai's Class 1 (``gentrifier proxy'') LQ index rose from 1.29 (1990) to 1.32 (2000), with Jing'an, Xuhui, Hongkou, and Luwan reaching LQ values as high as 1.45--1.56; the central-area Class 2 (``displacee proxy'') LQ fell from 0.95 to 0.79 over the same period, while outskirt districts' Class 2 LQ rose from 1.07 to 1.15 (Table 4).
  \item \textbf{Industrial upgrading}: Real estate management and consultancy employment grew by 145.29\% between 1990 and 2000 and by a further 417.80\% between 2000 and 2007; banking and insurance employment grew by 192.13\% (1990--2000); manufacturing, electricity, gas, and water production/supply employment fell by 24.34\% (1990--2000) and a further 6.59\% (2000--2007) (Table 5).
  \item \textbf{Liangwancheng vs.\ Shenjiazhai contrast}: 72.4\% of Liangwancheng survey respondents are in professional/managerial occupations versus 16\% in Shenjiazhai; 63.5\% of Liangwancheng respondents hold a college/university degree or above, versus a small minority in Shenjiazhai; 94.5\% of Liangwancheng housing is purchased commodity housing versus 70\% privately owned (mostly inherited) housing in Shenjiazhai; second-home ownership is 25.5\% in Liangwancheng versus 4\% in Shenjiazhai; by end-2007, Liangwancheng housing prices exceeded 23,000 Yuan per square metre, while the average estimated value of an entire inherited property in Shenjiazhai (approximately 250,000 Yuan) was comparable to the price of a single bathroom in Liangwancheng (Table 7).
  \item \textbf{Xintiandi vs.\ Jing'anli contrast}: 87\% of Xintiandi respondents are under 40 years old, versus only 19\% in Jing'anli (where 55\% are aged 40--65 and 26\% are over 65); 78\% of Xintiandi respondents hold a college/university degree or above (15\% at master's level or beyond) versus 13.1\% in Jing'anli; two-thirds of Xintiandi respondents report annual family income above 50,000 Yuan (35\% above 200,000 Yuan), versus a majority of Jing'anli respondents below 50,000 Yuan (Table 7). Because Xintiandi was converted into a purely commercial/recreational district, this case shows exclusionary displacement of the working-class user base rather than direct residential eviction at the flagship site itself, even though the wider Taipingqiao project produced substantial direct residential displacement.
  \item \textbf{Predicted impacts of displacement}: Among Jing'anli and Shenjiazhai residents, 34.1\% and 32.5\% respectively predicted inconvenient access to public facilities as the most prominent consequence of displacement, followed by risk of unemployment (24.8\%/28\%), fragmented social networks (22.9\%/23.5\%), and increased distance from the city center (18.2\%/16\%) (Table 8).
  \item \textbf{Qualitative evidence of livelihood harm}: In-depth interviews with residents relocated to peri-urban Pudong and Minhang document job loss due to unaffordable commuting costs and time, loss of access to cultural amenities (e.g., regular attendance at Yue opera), reliance on the Minimum Living Standard Scheme subsidy among relocated low-income households, and a residents'-committee officer's observation that relocation exacerbated unemployment because remote location and underdeveloped infrastructure foreclosed informal income-generating activities (e.g., small business, room renting).
\end{enumerate}

\subsection*{Contributions}

\begin{enumerate}
  \item \textbf{Empirical contribution}: Provides the first systematic census-based Location Quotient evidence of occupational class remaking across central Shanghai (1990--2000/2007), paired with micro-level household survey and interview evidence contrasting socioeconomic profiles across four matched pre- and post-gentrification neighborhoods.
  \item \textbf{Theoretical contribution}: Extends the ``new-build gentrification'' concept (Davidson and Lees, 2005) to a state-socialist transitional economy, arguing that Chinese gentrification is fundamentally \emph{state-sponsored} --- driven by municipal land supply, infrastructure investment, and resettlement housing provision --- rather than a straightforward import of Western market-driven or globalization-driven gentrification dynamics.
  \item \textbf{Methodological contribution}: Adapts the LQ-based occupational gentrifier/displacee proxy method (Atkinson, 2000) to a data-constrained transitional context that lacks direct statistical tracking of gentrification or displacement, and combines this with matched-neighborhood survey and interview data to substitute for the absence of longitudinal, individual-level panel data on displaced households.
  \item \textbf{Policy contribution}: Critiques Shanghai's ``365 scheme'' and ``new round old urban area redevelopment scheme'' as depoliticized, beautification-framed policy instruments that mask the socio-spatial polarization and welfare costs imposed on displaced low-income residents, and calls for renewed critical, class-conscious gentrification research (echoing Slater, 2006, 2008; Wacquant, 2008).
\end{enumerate}

\subsection*{Limitations and Critical Assessment}

\begin{enumerate}
  \item \textbf{Identification threats}: The study offers no causal identification strategy. The neighborhood pairs (Shenjiazhai/Liangwancheng; Jing'anli/Xintiandi) are matched only on pre-existing built environment and development history, not on other plausible confounders such as underlying land value, proximity to the Bund/CBD, or waterfront access. Because redevelopment sites are not randomly assigned but are selected by the state and developers precisely because of favorable location combined with ``dilapidated'' status, the observed socioeconomic contrasts may partly reflect pre-existing locational advantage and selective targeting of sites with the largest latent land-value gap, rather than purely the causal effect of redevelopment itself. The paper does not disentangle site-selection effects from treatment effects.
  \item \textbf{External validity}: The evidence is drawn from two redevelopment projects in a single city under a distinctive hybrid state-socialist/market-transition land tenure regime (state land ownership, leasehold conveyance, household registration/\emph{hukou}-linked entitlements, and top-down urban regeneration governance). The author's broader claims about Chinese new-build gentrification as a national phenomenon rely substantially on citations to her own earlier single-city Shanghai studies (He, 2007; He and Wu, 2007a, 2007b) rather than independent replication in other Chinese cities, so generalization even within China should be treated as a working hypothesis rather than an established finding.
  \item \textbf{Omitted mechanisms}: The paper documents compositional outcomes and self-reported anticipated hardship but does not formally model the underlying land-value mechanism. No rent-gap calculation (capitalized versus potential ground rent) is presented, no analysis is offered of local-government capture of land-lease conveyance revenue (the ``land finance'' channel well documented elsewhere in Chinese urban political economy), and no quantification is given of whether displacement is attributable to gentrification-induced land-value change per se, or to inadequate compensation relative to replacement housing cost in the state-mediated relocation process --- a distinction that matters greatly for policy diagnosis.
  \item \textbf{Data constraints}: The Xintiandi sample's incomparability (consumers/visitors rather than residents) is a genuine constraint acknowledged by the author but not fully resolved. Sample sizes per site (100--200 households) are small, drawn from single housing estates via non-random, purposive selection, which limits inference even within Shanghai. The Class 1/Class 2 occupational dichotomy is a coarse proxy for gentrifier/displacee status that does not directly observe residential mobility histories or migration trajectories. Critically, the reported ``impacts of displacement'' in Table 8 for Jing'anli and Shenjiazhai residents are \emph{anticipated}, hypothetical assessments by residents who have not yet been displaced, not measured post-relocation outcomes; only the smaller set of in-depth interviews captures actual post-relocation experience, and these are illustrative rather than representative. The 2000 occupational census figures are a rescaled 10\% sample, adding measurement noise to small-area LQ estimates.
  \item \textbf{Equilibrium considerations}: The paper does not analyze general-equilibrium or filtering responses in the wider metropolitan housing market. It does not examine whether or how outskirt housing prices, quality, or supply adjust as displaced low-income households relocate there (the rising outskirt Class 2 LQ in Table 4 is consistent with a filtering/sorting process but is not analyzed as one), nor does it consider metropolitan-scale land-price or housing-supply responses to concentrated central-city redevelopment, nor broader labor-market effects of relocating workers away from central-city employment clusters beyond the anecdotal commuting-cost evidence in the interviews.
\end{enumerate}

\subsection*{Cross-Disciplinary Bridges}

\begin{itemize}
  \item \textbf{Connection to Urban Economics}: He's account of the 1988 pioneering land-lease auction in Changning district and the 1992 opening of urban land to private developers, together with the ``re-emergence of a land price gradient'' in central Shanghai, parallels the bid-rent logic of the Alonso-Muth-Mills framework, though the paper never formalizes this as a bid-rent function or hedonic capitalization exercise. The shift from inherited/public rental tenure to purchased ``commodity housing'' documented in Table 7 is suggestive of a filtering-type tenure transition that a hedonic or structural housing-market model could test directly against the reported price data (4,252 Yuan/m\textsuperscript{2} in the 2004 survey rising to over 23,000 Yuan/m\textsuperscript{2} by 2007 in Liangwancheng).
  \item \textbf{Connection to Regional Science}: The Location Quotient methodology and sub-district-level choropleth mapping (Figure 2) are core regional-science spatial-concentration tools, and the concentric-ring population analysis around People's Square (2.5/5/7.5/10\,km) is a monocentric spatial-structure exercise. A regional-science extension would formalize spatial dependence with Moran's I or LISA statistics and a spatial weights matrix rather than relying on descriptive LQ maps alone.
  \item \textbf{Connection to Geographic Finance}: The land-lease-financed redevelopment model documented here is a direct empirical antecedent of the Chinese ``land finance'' literature on municipal reliance on land-conveyance revenue to fund urban infrastructure, and connects to work on state-mediated risk transfer to displaced households who receive resettlement housing/compensation rather than market-priced indemnification.
  \item \textbf{Suggested related works}:
    \begin{itemize}
      \item Wu, F. (2004). Residential relocation under market-oriented redevelopment: the process and outcomes in urban China. \emph{Geoforum}, 35 --- a direct empirical complement on the relocation and compensation process that He's interviews only sketch anecdotally.
      \item Smith, N. (1979). Toward a theory of gentrification. \emph{Antipode} --- the rent gap framework implicitly relevant to, but never formalized in, He's account of why ``dilapidated'' central sites are targeted for state-led redevelopment.
      \item Hsieh, C.-T. and Moretti, E. (2019). Housing constraints and spatial misallocation. \emph{American Economic Journal: Macroeconomics} --- a regional-science/urban-economics lens on the aggregate efficiency costs of restricting low-income access to high-productivity central-city locations, relevant to interpreting the welfare costs of the displacement documented here.
    \end{itemize}
\end{itemize}

\subsection*{Key Equations}

The paper's only formal quantitative construct is the Location Quotient (LQ) index (following Burt, 1996), used to measure the relative concentration of an occupational or industrial group in a given spatial unit compared with its concentration across the whole metropolitan area:

\[
LQ_{j,k} = \frac{P_{j,k} / P_{j}}{P_{k} / P}
\]

where $P_{j,k}$ is the population of occupational/industrial group $k$ in spatial unit $j$ (e.g., a district or sub-district), $P_{j}$ is the total population aged 15--65 in spatial unit $j$, $P_{k}$ is the population of group $k$ across the whole metropolitan area, and $P$ is the total metropolitan-area population aged 15--65. An $LQ_{j,k}$ value above 1 indicates that group $k$ is over-represented in unit $j$ relative to the metropolitan average; He applies this separately to Class 1 (gentrifier proxy) and Class 2 (displacee proxy) occupational groupings.

\subsection*{Quotable Passages}

\begin{quote}
"The city is celebrating its thriving neo-liberal urbanism by implementing enormous new-build gentrification, mostly in the form of demolition--rebuild development involving direct displacement of residents and landscapes." (p.~345)
\end{quote}

\begin{quote}
"New-build gentrification in China did not emerge in a post-regression era, but rather is the product of a rapidly developing society dominated by the imperative of capital accumulation and economic growth." (p.~358)
\end{quote}

\begin{quote}
"While the city celebrates its thriving neo-liberal urbanism, the long-term socioeconomic benefits of local communities are sacrificed." (p.~359)
\end{quote}
